\section{Duality of Feature and Sample Screening}
\label{sec:screening}

This section shows that feature screening and sample screening form an
equivariant pair under FR duality. We proceed in four steps. First, we
give a formal definition of feature and sample screening, independent
of any specific application problem. Second, we define feature and
sample translation and establish their equivariance. Third, we show
that both screenings decompose into a translation followed by a
restriction. Finally, combining these two facts, we prove the main
duality identity. As a by-product, we also discuss the resulting
reduction in intrinsic complexity and a characterization of
simultaneous screening.


\subsection{Feature and Sample Screening as Transformations}

This subsection defines feature screening and sample screening as
transformations on \(\Gamma\). Recall that an FR
representation \(p=(f,g,A,a)\) encodes the objective function
\[
p(x)=f(Ax)+g(x)+a.
\]

In the literature, feature screening typically eliminates a block of primal variables. We state the
assumption slightly more generally, allowing the eliminated block to
sit at an arbitrary fixed point rather than requiring it to be zero.
Let \(I\sqcup\bar I=\{1,\ldots,n\}\) and \(z\in\mathbb R^{\bar I}\) be a known vector, and
suppose that
\[
x_{\bar I}=z.
\]
Writing \(x=(x_I,x_{\bar I})\), \(A=(A_I,A_{\bar I})\), and
\(g(x)=g_I(x_I)+g_{\bar I}(x_{\bar I})\) by separability, substituting
\(x_{\bar I}=z\) into the objective gives
\[
p(x_I,z)
=
f(A_Ix_I+A_{\bar I}z)+g_I(x_I)+a+g_{\bar I}(z).
\]

\begin{definition}[Feature screening]
The \emph{feature screening} of \(p\) at \((I,z)\) is the FR
representation
\[
S^{F}_{I,z}(p)
:=
\bigl(
T^{\mathrm d}_{A_{\bar I}z}f,\;
g_I,\;
A_I,\;
a+g_{\bar I}(z)
\bigr).
\]
\end{definition}

Here, \(S^F_{I,z}(p)(x_I)=p(x_I,z)\). This means the screened objective is exactly the original objective with \(x_{\bar I}\) fixed at \(z\).

\begin{example}[Feature screening for the Lasso]
\label{ex:lasso-feature}
Classical feature screening for the Lasso is the special case \(z=0\)
of the proposed feature screening operator \(S^F_{I,z}\). The Lasso
problem
\[
\min_{x\in\mathbb R^n}
\frac12\|Ax-b\|_2^2+\lambda\|x\|_1,
\qquad \lambda>0,
\]
is an FR representation \(p=(f,g,A,0)\) with
\(f(y)=\frac12\|y-b\|_2^2\) and \(g(x)=\lambda\|x\|_1\).

In sparse learning, feature screening certifies that a block of
coefficients is inactive at the optimum. For the Lasso, inactivity means
that the corresponding coefficients vanish,
\[
x_{\bar I}=0,
\]
exactly the special case \(z=0\) of the general feature-screening
assumption \(x_{\bar I}=z\). Since \(z=0\), the screened function
\(S^F_{I,0}(p)\) is simply the restriction of the Lasso representation
to the remaining coordinates, with corresponding problem
\[
\min_{x_I\in\mathbb R^{|I|}}
\frac12\|A_Ix_I-b\|_2^2+\lambda\|x_I\|_1.
\]
The proposed definition thus recovers classical Lasso feature
screening, while extending it from coefficients certified to be zero
to coefficients certified to take any fixed value \(z\).
\end{example}


Feature screening simplifies the objective function by fixing a block of the variables in \(x\). Sample screening simplifies the objective function by fixing a block of the loss function \(f\). Let
\(J\sqcup\bar J=\{1,\ldots,m\}\) and \(s\in\mathbb R^{\bar J}\), and
suppose that
\[
s\in\partial f_{\bar J}(A_{\bar J}x),
\]
i.e., \(s\) is a fixed subgradient of the block \(\bar J\) of \(f\) evaluated at
\(x\). By the Fenchel--Young equality, which holds
in \(f(y)+f^*(s)\geq\langle s,y\rangle\) exactly when \(s\in\partial
f(y)\), the corresponding block of \(f\) can be replaced by a simple affine function, i.e.,
\[
f_{\bar J}(A_{\bar J}x)
=
\langle s,A_{\bar J}x\rangle-f_{\bar J}^{*}(s).
\]
Substituting this into the objective gives
\[
p(x)
=
f_J(A_Jx)+g(x)+\langle A_{\bar J}^{\top}s,x\rangle+a-f_{\bar J}^{*}(s).
\]


\begin{definition}[Sample screening]
The \emph{sample screening} of \(p\) at \((J,s)\) is the FR
representation
\[
S^{S}_{J,s}(p)
:=
\bigl(
f_J,\;
T^{\mathrm s}_{A_{\bar J}^{\top}s}g,\;
A_J,\;
a-f_{\bar J}^{*}(s)
\bigr),
\]
so that \(S^S_{J,s}(p)(x)\) equals the original objective with
\(f_{\bar J}(A_{\bar J}x)\) replaced by its Fenchel--Young-tight value
\(\langle s,A_{\bar J}x\rangle-f_{\bar J}^{*}(s)\).
\end{definition}

Geometrically, sample screening replaces the eliminated loss block by
its exact affine value at the certified dual point \(s\).

\begin{example}[Sample screening for the soft-margin SVM]
\label{ex:svm-sample}
Classical sample screening for the soft-margin SVM is the special case
\(s=0\) of the proposed sample screening operator \(S^S_{J,s}\). The
soft-margin SVM problem
\[
\min_{x\in\mathbb R^n}
\frac12\|x\|_2^2
+
C\sum_{i=1}^m
\max\!\left(0,\,1-b_i\langle a_i,x\rangle\right),
\qquad C>0,
\]
is represented in FR form by \(p=(f,g,A,0)\), where
\(f(y)=C\sum_{i=1}^m\max(0,1-b_i y_i)\), \(g(x)=\frac12\|x\|_2^2\), and
\(A=(a_1,\ldots,a_m)^\top\).

From \citep{ogawa2013}, the idea of sample screening is that, if one can identify a sample \(i\) that is correctly classified, i.e.,
\[1-b_i\langle a_i,x\rangle<0,\]
then one can remove the corresponding loss term from the loss function.

For \(J\sqcup\bar J=\{1,\ldots,m\}\), applying the above sample
screening condition for all indices \(i\in\bar J\) can be rewritten as the following condition\footnote{In the original paper~\citep{ogawa2013}, the authors also screen samples such that \(1\in\partial f_{\bar J}(A_{\bar J}x)\). Our sample screening operator also covers this case; for simplicity, we do not discuss it in this example.}
\[
0\in\partial f_{\bar J}(A_{\bar J}x).
\]
This is exactly the special case \(s=0\) of the generalized assumption
\(s\in\partial f_{\bar J}(A_{\bar J}x)\). In this case, the translation
terms vanish, \(A_{\bar J}^{\top}s=0\) and \(f_{\bar J}^*(s)=0\), so
\(S^S_{J,0}(p)=(f_J,g,A_J,0)\), and the screened problem is
\[
\min_{x\in\mathbb R^n}
\frac12\|x\|_2^2
+
C\sum_{i\in J}
\max\!\left(0,\,1-b_i\langle a_i,x\rangle\right).
\]
The proposed definition thus recovers classical SVM sample screening,
while extending it from samples certified to satisfy \(0\in\partial
f_i\) to samples certified to satisfy an arbitrary subgradient
condition \(s\in\partial f_i\).
\end{example}

Both definitions already display the matrix slot explicitly, so the
effect of screening on problem size needs no further argument.

\begin{corollary}[Screening reduces problem size]
\label{cor:complexity}
For every \(p=(f,g,A,a)\in\Gamma_{m,n}\), every
\(I\sqcup\bar I=\{1,\ldots,n\}\), and every \(J\sqcup\bar J=\{1,\ldots,m\}\),
\[
\operatorname{size}\bigl(S^F_{I,z}(p)\bigr)=(m,|I|),
\qquad
\operatorname{size}\bigl(S^S_{J,s}(p)\bigr)=(|J|,n).
\]
\end{corollary}

\begin{proof}
Immediate from the definitions above: \(S^F_{I,z}(p)\) has matrix slot
\(A_I\in\mathbb R^{m\times|I|}\), and \(S^S_{J,s}(p)\) has matrix slot
\(A_J\in\mathbb R^{|J|\times n}\).
\end{proof}
\noteCLAUDE{Doc2 dependencies: none. Uses only the definitions of
\(S^F_{I,z}\), \(S^S_{J,s}\) from this section.}

Reading intrinsic complexity as problem size, screening rules reduce
the intrinsic complexity. Theorem~\ref{thm:decomposition} below
explains why this must always happen, since translation never changes
the shape of \(A\) and restriction is the only step that can, but the
size claim itself needs nothing beyond the definitions just given.

Both definitions arose from a direct assumption on \(p\), a primal
value for feature screening and a dual subgradient for sample
screening, with no reference to translation or restriction. Yet both
formulas visibly split into "shift \(f\) or \(g\), then drop
coordinates." Section~\ref{sec:operators} introduced exactly two
primitives for this: translation (\(T^{\mathrm d}\), \(T^{\mathrm v}\),
\(T^{\mathrm s}\), combined into \(T_k\)) and restriction (\(R^F\),
\(R^S\)). Theorem~\ref{thm:decomposition} shows that \(S^F_{I,z}\) and
\(S^S_{J,s}\), defined above purely from optimization assumptions, are
always built from exactly these two primitives and nothing else.



\subsection{Feature--Sample Translation}
In Subsection~\ref{subsec:translation}, we introduced the translation \(T_k\) and its translation parameter \(k\). In this subsection, we introduce feature and sample translation, built on top of \(T_k\) and \(k\).
Surprisingly, these objects also admit equivariance under FR duality.
These notions play an important role in establishing the main results proved in the following subsections.

\begin{definition}[Feature--Sample Translation]
    \label{def:feature-sample-translation}
    We define the feature translation
    \begin{equation*}
        T^F_{I, z}(p) = T_k(p), \qquad
        k = \kappa^F_{I, z}(p) :=  (A_{\bar I}z,\;g_{\bar I}(z),\;0)\in K_{m, |\bar I|}.
    \end{equation*}
    where \(k\) is a feature translation parameter depending on \(p\), index set \(I\) and \(z\in \mathbb R^{\bar I}\).

    Symmetrically, we define the sample translation
    \begin{equation*}
        T^S_{J, s}(p) = T_h(p), \qquad
        h = \kappa^S_{J, s}(p) := (0,\;-f_{\bar J}^{*}(s),\;A_{\bar J}^{\top}s)\in K_{|\bar J|, n}.
    \end{equation*}
    where \(h\) is a sample translation parameter depending on \(p\), index set \(J\) and \(s\in \mathbb R^{\bar J}\).
\end{definition}


Now we show that there is an equivariance structure on both the
parameter level and the translation level.

\begin{proposition}[Parameter Equivariance]
\label{prop:parameter-equivariance}
For every \(p=(f,g,A,a)\in\Gamma_{m,n}\), every
\(I\sqcup\bar I=\{1,\ldots,n\}\), \(z\in\mathbb R^{\bar I}\), every
\(J\sqcup\bar J=\{1,\ldots,m\}\), and \(s\in\mathbb R^{\bar J}\),
\begin{align*}
        (\kappa^F_{I, z}(p))^* = \kappa^S_{I, z}(p^*),\\
        (\kappa^S_{J, s}(p))^* = \kappa^F_{J, s}(p^*).
\end{align*}
\end{proposition}

\begin{proof}
Using \(k^*=(-d,-c,-b)\) for \(k=(b,c,d)\) (Section~\ref{sec:operators}),
\[
(\kappa^F_{I,z}(p))^*
=
(A_{\bar I}z,\,g_{\bar I}(z),\,0)^*
=
(0,\,-g_{\bar I}(z),\,-A_{\bar I}z).
\]
Write \(p^*=(g^*,f^*,-A^\top,-a)\). By definition, the sample
translation parameter of \(p^*\) at \((I,z)\) is
\[
\kappa^S_{I,z}(p^*)
=
\bigl(0,\;-(g^*)_{\bar I}^{*}(z),\;(-A^\top)_{\bar I}^{\top}z\bigr).
\]
Since \(g\) is separable, its conjugate restricts termwise,
\((g^*)_{\bar I}=(g_{\bar I})^*\), so
\((g^*)_{\bar I}^{*}(z)=(g_{\bar I})^{**}(z)=g_{\bar I}(z)\) by
biconjugation, using that \(g_{\bar I}\), a finite sum of closed,
proper, convex functions, is itself closed, proper, convex\noteLE{Basic
convex-analysis fact, citation to follow.}. Rows
\(\bar I\) of \(-A^\top\) equal \(-(A_{\bar I})^\top\), so
\((-A^\top)_{\bar I}^\top z=-A_{\bar I}z\). Hence
\[
\kappa^S_{I,z}(p^*)=(0,\,-g_{\bar I}(z),\,-A_{\bar I}z)=(\kappa^F_{I,z}(p))^*.
\]

We now prove the second identity, which does not follow from the first
by simply exchanging \(f\) and \(g\): Definition~\ref{def:feature-sample-translation}
already builds \(\kappa^F_{I,z}\) and \(\kappa^S_{J,s}\) asymmetrically,
so it needs its own computation.
Using \(k^*=(-d,-c,-b)\) again,
\[
(\kappa^S_{J,s}(p))^*
=
(0,\,-f_{\bar J}^{*}(s),\,A_{\bar J}^{\top}s)^*
=
(-A_{\bar J}^{\top}s,\,f_{\bar J}^{*}(s),\,0).
\]
Write \(p^*=(g^*,f^*,-A^\top,-a)\in\Gamma_{n,m}\). Since \(J\subseteq
\{1,\ldots,m\}\) indexes the second slot's problem size for \(p^*\), the
feature translation parameter of \(p^*\) at \((J,s)\) is
\[
\kappa^F_{J,s}(p^*)
=
\bigl((-A^\top)_{\bar J}s,\;(f^*)_{\bar J}(s),\;0\bigr),
\]
by the same definition as \(\kappa^F_{I,z}\), with \(g^*\) taking the
role of the first slot's function and \(-A^\top\) the role of the
matrix. Columns \(\bar J\) of \(-A^\top\) equal \(-(A_{\bar J})^\top\),
since columns of \(A^\top\) are rows of \(A\) transposed, so
\((-A^\top)_{\bar J}s=-A_{\bar J}^\top s\). Since \(f\) is separable, its
conjugate restricts termwise, \((f^*)_{\bar J}=(f_{\bar J})^*\), so
\((f^*)_{\bar J}(s)=f_{\bar J}^{*}(s)\); this is the same separability
fact used above, now applied to \(f\) rather than \(g\), and used
directly rather than through biconjugation, since here it is the
conjugate \(f^*\) itself, not its biconjugate, that must be restricted.
Hence
\[
\kappa^F_{J,s}(p^*)=(-A_{\bar J}^{\top}s,\,f_{\bar J}^{*}(s),\,0)=(\kappa^S_{J,s}(p))^*.
\]
Unlike the first identity, this computation never invokes
biconjugation: \(\kappa^S_{J,s}(p)\) already carries a conjugate,
\(f_{\bar J}^*(s)\), in its second slot, so dualizing and restricting
commute with a single separability step, not two.
\end{proof}

Proposition~\ref{prop:parameter-equivariance} identifies where the
feature--sample asymmetry lives. Feature screening produces a domain
shift and a value shift, while sample screening produces a slope shift
and a conjugate value shift. The proposition shows that this apparent
asymmetry is fully absorbed by the canonical translation parameters:
FR duality exchanges the feature parameter with the sample parameter.
The next step is to lift this parameter equivariance from parameters to
translation operators.

\begin{proposition}[Feature--Sample Translation Equivariance]
\label{prop:translation-equivariance-screening}
For every \(p=(f,g,A,a)\in\Gamma_{m,n}\), every
\(I\sqcup\bar I=\{1,\ldots,n\}\), \(z\in\mathbb R^{\bar I}\), every
\(J\sqcup\bar J=\{1,\ldots,m\}\), and \(s\in\mathbb R^{\bar J}\),
\begin{align*}
        (T^F_{I, z}(p))^* = T^S_{I, z}(p^*),\\
        (T^S_{J, s}(p))^* = T^F_{J, s}(p^*).
\end{align*}
\end{proposition}

\begin{proof}
By Translation Equivariance (Proposition~\ref{thm:fenchel-equivariance},
Section~\ref{sec:operators}) applied to \(k=\kappa^F_{I,z}(p)\),
\[
(T^F_{I,z}(p))^*
=
\bigl(T_{\kappa^F_{I,z}(p)}(p)\bigr)^*
=
T_{(\kappa^F_{I,z}(p))^*}(p^*).
\]
By Parameter Equivariance (Proposition~\ref{prop:parameter-equivariance}),
\((\kappa^F_{I,z}(p))^*=\kappa^S_{I,z}(p^*)\), so
\[
(T^F_{I,z}(p))^*
=
T_{\kappa^S_{I,z}(p^*)}(p^*)
=
T^S_{I,z}(p^*).
\]
The second identity follows the same argument, exchanging the roles of
\(f\) and \(g\), \(I\) and \(J\), and domain and slope translation.
\end{proof}




\subsection{Decomposition of Screening}

This subsection proves that feature screening and sample screening are
not primitive: each decomposes into a translation followed by a
restriction. 


\begin{theorem}[Screening Decomposition]
\label{thm:decomposition}
For every \(p=(f,g,A,a)\), every \(I\sqcup\bar I=\{1,\ldots,n\}\),
\(z\in\mathbb R^{\bar I}\), \(J\sqcup\bar J=\{1,\ldots,m\}\), and
\(s\in\mathbb R^{\bar J}\),
\[
S^{F}_{I,z}
=
R^{F}_{I}\circ T^{F}_{I,z},
\qquad
S^{S}_{J,s}
=
R^{S}_{J}\circ T^{S}_{J,s}.
\]
\end{theorem}

\begin{proof}
By Definition~\ref{def:feature-sample-translation}, \(\kappa^F_{I,z}(p)\) has zero slope component, so
\[
T^F_{I,z}(p)
=
\bigl(T^{\mathrm d}_{A_{\bar I}z}f,\;g,\;A,\;a+g_{\bar I}(z)\bigr).
\]
Applying \(R^F_I\), which restricts the second and third slots to
\(I\) and leaves the first and fourth untouched,
\[
R^F_I\bigl(T^F_{I,z}(p)\bigr)
=
\bigl(T^{\mathrm d}_{A_{\bar I}z}f,\;g_I,\;A_I,\;a+g_{\bar I}(z)\bigr),
\]
which is exactly \(S^F_{I,z}(p)\) from the definition above.

The sample case is the mirror computation. Since \(\kappa^S_{J,s}(p)\)
has zero domain component,
\[
T^S_{J,s}(p)
=
\bigl(f,\;T^{\mathrm s}_{A_{\bar J}^{\top}s}g,\;A,\;a-f_{\bar J}^{*}(s)\bigr).
\]
Applying \(R^S_J\), which restricts the first and third slots to \(J\)
and leaves the second and fourth untouched,
\[
R^S_J\bigl(T^S_{J,s}(p)\bigr)
=
\bigl(f_J,\;T^{\mathrm s}_{A_{\bar J}^{\top}s}g,\;A_J,\;a-f_{\bar J}^{*}(s)\bigr),
\]
which is exactly \(S^S_{J,s}(p)\).
\end{proof}
\noteCLAUDE{Doc2 dependencies: Definition "Translation" (\(T_k(p)\))
and Definition "Feature and sample restriction" (\(R^F_I\), \(R^S_J\))
only. No proposition or lemma needed.} \noteLE{Verified.}

This is the paper's central structural fact, screening equals
translation followed by restriction, now proved rather than assumed.
Feature screening and sample screening were defined from two
unrelated-looking optimization assumptions, a primal fixed point and a
dual subgradient, and Theorem~\ref{thm:decomposition} shows they are
the same two-step construction, one primitive translation and one
primitive restriction, applied to the two different primitives of
Section~\ref{sec:operators}. This is what lets the rest of this
section's results, Feature--Sample Duality and Simultaneous Screening
Equivariance, follow from Section~\ref{sec:operators}'s Translation
Equivariance and Restriction Equivariance, rather than needing a
separate argument for each screening type.

Theorem~\ref{thm:decomposition} also gives a second proof of
Corollary~\ref{cor:complexity}, stated already above directly from the
definitions: since translation never changes the shape of \(A\)
(Section~\ref{sec:operators}), the decomposition shows the size change
comes entirely from the restriction half, \(R^F_I\) or \(R^S_J\),
consistent with, and not needed beyond, the direct argument already
given.



\subsection{Duality of Screening}

This subsection proves that feature screening and sample screening are
equivariant under FR duality: dualizing one produces exactly the other,
on the dual representation. Feature screening moves the eliminated
block of \(x\) into a domain translation of \(f\), then discards it
from \(g\). FR duality exchanges domain translations
with slope translations and exchanges the roles of \(f\) and \(g\)
(Section~\ref{sec:operators}). Sample screening moves an eliminated
block of \(y\) into a slope translation of \(g\), then discards it
from \(f\). Matching these two descriptions suggests that dualizing a
feature screening produces exactly a sample screening, with no other
operation involved.

\begin{theorem}[Feature--Sample Duality]
\label{thm:duality}
Let \(p=(f,g,A,a)\in\Gamma_{m,n}\). For every
\(I\sqcup\bar I=\{1,\ldots,n\}\) and \(z\in\mathbb R^{\bar I}\),
\[
\left(S^{F}_{I,z}(p)\right)^*
=
S^{S}_{I,z}(p^*).
\]
Symmetrically, for every \(J\sqcup\bar J=\{1,\ldots,m\}\) and
\(s\in\mathbb R^{\bar J}\),
\[
\left(S^{S}_{J,s}(p)\right)^*
=
S^{F}_{J,s}(p^*).
\]
\end{theorem}

\begin{proof}
We prove the first identity; the second is the mirror computation with
the roles of \(f\) and \(g\), \(I\) and \(J\), domain and slope
translation, exchanged. The proof combines Screening Decomposition
(Theorem~\ref{thm:decomposition}, this section) with Restriction
Equivariance (Section~\ref{sec:operators}) and Feature--Sample
Translation Equivariance
(Proposition~\ref{prop:translation-equivariance-screening}, this
section).

By Theorem~\ref{thm:decomposition}, \(S^F_{I,z}(p)=R^F_I(T^F_{I,z}(p))\).
By Restriction Equivariance, applied to \(q:=T^F_{I,z}(p)\),
\[
\left(S^F_{I,z}(p)\right)^*
=
\left(R^F_I(q)\right)^*
=
R^S_I(q^*)
=
R^S_I\bigl((T^F_{I,z}(p))^*\bigr).
\]
By Feature--Sample Translation Equivariance,
\((T^F_{I,z}(p))^*=T^S_{I,z}(p^*)\). Substituting,
\[
\left(S^F_{I,z}(p)\right)^*
=
R^S_I\bigl(T^S_{I,z}(p^*)\bigr)
=
S^S_{I,z}(p^*),
\]
the last equality by Theorem~\ref{thm:decomposition} applied to \(p^*\).
This proves the first identity.

The second identity is the mirror computation, exchanging the roles of
\(f\) and \(g\), \(I\) and \(J\), and domain and slope translation,
using the second halves of Screening Decomposition, Restriction
Equivariance, and Feature--Sample Translation Equivariance.
\end{proof}
\noteCLAUDE{Dependencies: Theorem~\ref{thm:decomposition} and
Proposition~\ref{prop:translation-equivariance-screening} (this
section), plus Restriction Equivariance (doc2). No convex-analysis fact
is needed inline anymore: the separability/biconjugation argument now
lives entirely inside Proposition~\ref{prop:parameter-equivariance}'s
proof, cited transitively through
Proposition~\ref{prop:translation-equivariance-screening}.}


Theorem~\ref{thm:duality} says a single screening dualizes to a single
screening of the other type. What happens when a feature screening and a sample screening are both applied to the same \(p\)? 
This composition was used in \citep{shibagaki2016simultaneous} during the solving process, where it is referred to as simultaneous screening. The following result shows that, by exchanging their order with the appropriate parameters, the two compositions also form an equivariant pair.

\begin{corollary}[Simultaneous Screening Equivariance]
\label{cor:simultaneous}
For every \(p=(f,g,A,a)\in\Gamma_{m,n}\), every
\(I\sqcup\bar I=\{1,\ldots,n\}\), \(z\in\mathbb R^{\bar I}\),
\(J\sqcup\bar J=\{1,\ldots,m\}\), and \(s\in\mathbb R^{\bar J}\),
whenever both compositions below are defined,
\[
\left(S^{S}_{J,s}\,S^{F}_{I,z}(p)\right)^*
=
S^{F}_{J,s}\,S^{S}_{I,z}(p^*),
\]
equivalently,
\[
S^{S}_{J,s}\,S^{F}_{I,z}(p)
=
\left(S^{F}_{J,s}\,S^{S}_{I,z}(p^*)\right)^*.
\]
\end{corollary}

\begin{proof}
Write \(q:=S^F_{I,z}(p)\). Applying Theorem~\ref{thm:duality} to \(q\) and \(p\), respectively, we obtain
\[
\left(S^S_{J,s}(q)\right)^*=S^F_{J,s}(q^*),
\qquad 
q^*=\left(S^F_{I,z}(p)\right)^*=S^S_{I,z}(p^*).
\]
Substituting,
\[
\left(S^S_{J,s}\,S^F_{I,z}(p)\right)^*
=
S^F_{J,s}\bigl(S^S_{I,z}(p^*)\bigr)
=
S^F_{J,s}\,S^S_{I,z}(p^*),
\]
which is the first identity. The second follows by applying
\((\cdot)^*\) to both sides and using the involution \((p^*)^*=p\)
(Proposition~\ref{prop:involution}).
\end{proof}
% \noteCLAUDE{Doc2 dependencies: only Proposition "Involution"
% (\((p^*)^*=p\)). Everything else comes from Theorem~\ref{thm:duality}
% (this section), applied twice.}\noteLE{Verified}

% This is the correct replacement for a Commutativity theorem. Feature
% screening and sample screening do not need to commute on the same
% representation \(p\): requiring that would identify two representations
% that, in general, differ by exactly the cross-block interaction
% between the dropped samples and the dropped features. Instead,
% composing on \(p\) and dualizing lands on the same object as composing
% on \(p^*\), with \(S^F\) and \(S^S\) exchanging roles. This is what
% "simultaneous screening" means in this calculus: not that the two
% operators commute, but that composing them is itself compatible with
% duality, exactly as a single screening is (Theorem~\ref{thm:duality}).
% This is a two-line corollary of Theorem~\ref{thm:duality}, applied
% twice; no new machinery is needed beyond it.
